% !TeX root = 20260825.tex
% !LW recipe=lablog XeLaTeX
\input{_support/style.tex}

\begin{document}

\reporttitle{2026-08-25}{실험 결과 및 설계 기록}

\section*{목적과 이 보고서의 범위}

이 날짜의 실험 설정, 관찰 결과, 원인에 대한 해석과 후속 판단을 함께 정리한다. 다른 날짜의 평가 방식이나 표본 수가 다른 점수를 직접 합치지 않으며, 실행 완료와 논문 근거 채택을 구별한다.

\recordnote{수합·문서 개편: 2026-10-20. 실행 날짜는 각 절과 원시 실행 목록을 따른다. 아래의 과거 계획·진행 상태는 당시 기록이며, 현재 채택 상태와 최종 결과는 후속 실험을 함께 확인한다.}

\section{2026-08-25 · 2026-08-25}
\label{report:20260825_0825}
\recordnote{이 절은 당시 설정의 결과다. 평가 방식 변경·정정·후속 비교는 뒤에 명시한 실험과 구별한다.}
\subsection{§1. 연구방향 PDF (docs/\allowbreak{}연구방향.pdf, RACE-PFN) 독해 Q\&A}
사용자 질문 4건 정리 (구현 전 합의 사항):

\par\noindent \textbullet\enspace \textbf{"residual representation"의 의미 (Step 2)}: 클래스로 클러스터를 만드는 것이 아님. r\_\allowbreak{}i(=\allowbreak{} class-balanced weight × global CE loss, D\_\allowbreak{}mine에서 측정)는 "global이 이 sample을 얼마나 어렵게/\allowbreak{}중요하게 틀렸나"의 스칼라이고, z\_\allowbreak{}i =\allowbreak{} φ(x\_\allowbreak{}i)는 라벨과 무관한 \textbf{feature 공간 임베딩}(RobustScaler+\allowbreak{}PCA). 클러스터링은 z-공간에서 \textbf{r\_\allowbreak{}i-가중 k-means} — 즉 "global의 가중 loss 질량이 몰려 있는 feature 영역" K개를 찾는 것. 한 클러스터에 여러 클래스가 섞일 수 있음(예: temporally shifted benign/\allowbreak{}attack boundary). 데이터: z·클러스터링은 D\_\allowbreak{}mine 행, global은 D\_\allowbreak{}context로 구성, 라벨은 r\_\allowbreak{}i 가중치로만 개입.


\par\noindent \textbullet\enspace \textbf{Step 4는 학습이 아니라 라벨 제작}: D\_\allowbreak{}router에서 global과 모든 expert를 offline 실행해 Δ\_\allowbreak{}ik =\allowbreak{} L\_\allowbreak{}i0 − L\_\allowbreak{}ik − λc\_\allowbreak{}k를 계산 — 이것이 router(Step 5)와 verifier(Step 7)의 \textbf{학습 타깃}이 된다.


\par\noindent \textbullet\enspace \textbf{D\_\allowbreak{}router는 Step 5와 Step 7 양쪽에서 사용}: router(분류: NO/\allowbreak{}1..K, 타깃 k*=\allowbreak{}argmax Δ if max>0)와 verifier(회귀: Δ) 모두 D\_\allowbreak{}router로 학습. D\_\allowbreak{}context  \ensuremath{\rightarrow}  context 구성, D\_\allowbreak{}mine  \ensuremath{\rightarrow}  residual 채굴+\allowbreak{}expert block, D\_\allowbreak{}val  \ensuremath{\rightarrow}  K/\allowbreak{}예산/\allowbreak{}λ 선택, D\_\allowbreak{}test  \ensuremath{\rightarrow}  최종 1회 평가.


\par\noindent \textbullet\enspace \textbf{Step 7이 Huber/\allowbreak{}ranking인 이유}: verifier는 "예측이 맞나"(분류)가 아니라 \textbf{연속량 Δ(global 대비 loss 감소 기대치)를 회귀}한다. balanced CE 차는 꼬리가 무거워(rare-class −log p 폭발) L2가 outlier에 끌리므로 Huber로 강건화; 결정에 실제로 쓰이는 것은 Δ의 부호/\allowbreak{}순서라서 pairwise ranking도 대안. Δ̂>0 임계값은 "예상 loss 감소가 양수"라는 정의상 기준이지 튜닝 임계값이 아님(Step 8).


\subsection{§2. EXP22 사전 등록 — RACE-PFN v1 (결과 전 고정)}
사용자 지시: "docs/\allowbreak{}연구방향.pdf를 읽고 이 모델을 코드로 구현해줘." 스크립트: \texttt{tabpfn/\allowbreak{}\allowbreak{}n\allowbreak{}fv3\_\allowbreak{}\allowbreak{}v3\_\allowbreak{}\allowbreak{}e\allowbreak{}xp22\_\allowbreak{}\allowbreak{}rac\allowbreak{}epfn.\allowbreak{}py} (exp21 계보, TabPFN 사용 관용구·uniq\_\allowbreak{}rows· full\_\allowbreak{}proba 복사; 기존 파일 무수정). PDF §26 Phase 1–5 범위 구현: 자동 residual expert +\allowbreak{} counterfactual advantage routing +\allowbreak{} sparse top-1 +\allowbreak{} verifier. drift/\allowbreak{}OOD/\allowbreak{} context-refresh(Phase 6–7)는 후속.

\textbf{데이터 지도} (train pool 60\% 슬라이스를 클래스별 시간순 재분할; val/\allowbreak{}test 불변):

\par\noindent \textbullet\enspace D\_\allowbreak{}context 0.50  \ensuremath{\rightarrow}  global C0(예산 1M, stratified) · anchor A(클래스당 2k) · φ(RobustScaler+\allowbreak{}PCA16) fit · a0 affinity 참조행


\par\noindent \textbullet\enspace D\_\allowbreak{}mine 0.25  \ensuremath{\rightarrow}  r\_\allowbreak{}i =\allowbreak{} w\_\allowbreak{}y·CE(p0), w\_\allowbreak{}c=\allowbreak{}(N/\allowbreak{}(C·n\_\allowbreak{}c))\textasciicircum{}γ(γ=\allowbreak{}1, full-pool counts)  \ensuremath{\rightarrow}  r-가중 KMeans(K=\allowbreak{}4)  \ensuremath{\rightarrow}  S\_\allowbreak{}k =\allowbreak{} TopB(r\_\allowbreak{}i·exp(−d²/\allowbreak{}τ)), B=\allowbreak{}186k  \ensuremath{\rightarrow}  C\_\allowbreak{}k =\allowbreak{} A∪S\_\allowbreak{}k


\par\noindent \textbullet\enspace D\_\allowbreak{}router 나머지  \ensuremath{\rightarrow}  Δ\_\allowbreak{}ik 계산(λ\_\allowbreak{}cost=\allowbreak{}0 v1)  \ensuremath{\rightarrow}  router(XGB, h\_\allowbreak{}pre 30d) +\allowbreak{} shared verifier(XGB pseudo-Huber, h\_\allowbreak{}post 36d) — 시간순 80/\allowbreak{}20 fit/\allowbreak{}holdout


\par\noindent \textbullet\enspace D\_\allowbreak{}test  \ensuremath{\rightarrow}  global  \ensuremath{\rightarrow}  proposal  \ensuremath{\rightarrow}  top-1 expert만 실행  \ensuremath{\rightarrow}  Δ̂>0일 때만 override


\textbf{v1 단순화(은닉 아님, 문서화된 knob)}: a0/\allowbreak{}a\_\allowbreak{}k는 z-공간 16-NN 평균거리 proxy (TabPFN attention affinity는 후속 knob); λ\_\allowbreak{}cost=\allowbreak{}0(3a 아티팩트에 Δ 분포 기록 후 후속 run에서 설정); Huber(ranking은 ablation); context-selection=\allowbreak{}random(medoid 구현됨, 후속 원노브).

\textbf{판정 기준(고정)}: 같은-run 3중 튜플 — racepfn\_\allowbreak{}system vs global\_\allowbreak{}tabpfn vs full-pool XGBoost(같은 split). 관심: ① residual 클러스터가 수동 Flood/\allowbreak{}BF-Bot/\allowbreak{}Tail 구조를 재발견하는가(1b), ② positive-advantage 비율과 Δ 분포 (3a/\allowbreak{}3b), ③ router의 positive-advantage recall(4a), ④ verifier sign-accuracy(4b), ⑤ helpful/\allowbreak{}harmful override·net correction·ρ(5a/\allowbreak{}6a), ⑥ --dense-eval 시 oracle gap recovery. exp21 0.7411과는 \textbf{방향 비교만}(global context가 D\_\allowbreak{}context 반쪽 pool로 제한되어 수치 비교 불가 — split 사용법이 다름).

\textbf{알려진 실패 양식 대조(SRC\_\allowbreak{}HISTORY)}: \#1 routing 오류의 imbalance 증폭  \ensuremath{\rightarrow}  global-default +\allowbreak{} verifier 이중 안전판(override는 Δ̂>0에서만); \#2 val-예측 적합의 비전이  \ensuremath{\rightarrow}  router/\allowbreak{}verifier는 val이 아니라 \textbf{별도 D\_\allowbreak{}router}(context와 분리)로 학습, 시간순 holdout으로 내부 검증; \#6 long-tail energy gate 붕괴  \ensuremath{\rightarrow}  energy gate 자체를 제거하고 advantage 회귀로 대체.

스모크: cic2018\_\allowbreak{}capped 축소 예산 +\allowbreak{} --dense-eval로 전 단계 관통 확인 후 기록 run.

\subsection{§3. EXP22 구현 리뷰 +\allowbreak{} 스모크 결과 (인용 금지 — 기록용 아님)}
\textbf{적대적 리뷰} (14-agent workflow: 4개 렌즈 spec/\allowbreak{}leakage/\allowbreak{}runtime/\allowbreak{}conventions +\allowbreak{} 발견별 반박 검증): 확정 20건(반박 0), 중복 제거 9건 전부 수정 — ① router 타깃 비연속/\allowbreak{}2-클래스 시 XGBClassifier 크래시(xgb 3.2.0 실측 재현; dense encode +\allowbreak{} \texttt{present[\allowbreak{}]} 디코드 +\allowbreak{} 명시적 num\_\allowbreak{}class; 2-클래스 softprob의 predict가 float64를 반환하는 것도 실측으로 잡음), ② anchor가 클래스당 균등이 아니라 pool 비율 배분이던 것(IR>1000에서 anchor 취지 붕괴 — \texttt{anchor\_\allowbreak{}\allowbreak{}s\allowbreak{}ubset} 클래스당 min(2000, n\_\allowbreak{}c)로 교정, 2a에 anchor 조성 기록), ③ 무동작 --train-split/\allowbreak{} --skip-tabpfn이 Args 기록을 오염 가능  \ensuremath{\rightarrow}  SystemExit 가드, ④ seed 밴드 이중 적용· 충돌(밴드 100 간격 재배치, cap\_\allowbreak{}per\_\allowbreak{}class 내부 +\allowbreak{}500+\allowbreak{}i 고려), ⑤ h\_\allowbreak{}post의 expert-id one-hot을 문서화(공유 verifier의 per-expert offset 보정 — 의도적 스펙 추가), ⑥ 죽은 변수 제거, ⑦ report\_\allowbreak{}memory\_\allowbreak{}plan +\allowbreak{} expandable\_\allowbreak{}segments 경고 추가, ⑧ kmeans/\allowbreak{}medoid n\_\allowbreak{}init 등 argparse 노출, ⑨ propose() 디코딩 정정.

\textbf{스모크} (run \texttt{20260825\allowbreak{}\_\allowbreak{}\allowbreak{}133516\_\allowbreak{}\allowbreak{}nfv3\_\allowbreak{}\allowbreak{}cic\allowbreak{}2018capp\allowbreak{}ed\_\allowbreak{}\allowbreak{}exp22\allowbreak{}\_\allowbreak{}\allowbreak{}racepfn\allowbreak{}}; capped, C0 20k·anchor 300/\allowbreak{}class·block 4k·mine 60k·router 3k/\allowbreak{}class·eval 2k/\allowbreak{}class·dense):

\par\noindent \textbullet\enspace macro F1: \textbf{racepfn 0.8674} vs 같은-run global\_\allowbreak{}tabpfn 0.7710 (+\allowbreak{}0.0964) vs 같은-run XGB(full pool 361k) 0.8687 (−0.0013) vs dense oracle 0.9865. \textbf{oracle gap recovery 0.4473}. 클래스: benign .531 \ensuremath{\rightarrow} .705 · brute .670 \ensuremath{\rightarrow} .858 · dos .566 \ensuremath{\rightarrow} .795 · web .851 \ensuremath{\rightarrow} .944 · inf .781 \ensuremath{\rightarrow} .772(소폭 희생).


\par\noindent \textbullet\enspace routing: helpful 1,261 vs harmful 77 (\textbf{override precision 0.9425}, net +\allowbreak{}1,184); ρ=\allowbreak{}0.544 제안, verifier 통과 후 실행률 0.206; eval positive-advantage proposal recall 0.9678.


\par\noindent \textbullet\enspace residual 클러스터: 전-클래스 혼합 대형 1개(94\% 행, 질량 0.86) +\allowbreak{} 고잔차 소형 3개(dos/\allowbreak{}benign 경계·inf·web 포함) — capped라 IR이 낮아 uncapped에서 재확인 필요.


\textbf{주시 사항(기록 run에서 확인)}: ① verifier holdout \textbf{sign\_\allowbreak{}accuracy \textasciitilde{}0.17} (pearson r 0.36–0.75인데 부호는 반전 수준) — 시간순 fit/\allowbreak{}holdout 간 Δ 분포 이동의 신호로 보임; 시스템 지표(override precision 0.94)는 건재하나 verifier 보정이 chrono shift에 취약할 가능성. ② router holdout 제안률 0.0207 vs eval ρ 0.544 괴리 — 같은 원인 계열로 추정. 둘 다 D\_\allowbreak{}router의 시간 꼬리가 NO\_\allowbreak{}EXPERT 지배 구간인지 3a/\allowbreak{}3b로 확인할 것.

기록 run 발사: uncapped cic2018, 기본 예산(C0 1M · anchor 2k/\allowbreak{}class · block 186k · mine 2M · router 100k/\allowbreak{}class · eval full 4.02M) +\allowbreak{} --dense-eval (Phase 1 oracle 상보성 검증 포함).

\subsection{§4. EXP22 기록 run 결과 (run \texttt{20260825\allowbreak{}\_\allowbreak{}\allowbreak{}142540}, uncapped cic2018, test 4.02M) — \textbf{첫 구성 패배, oracle headroom 확인}}
\textbf{판정}: racepfn\_\allowbreak{}system macro \textbf{0.6346} — 같은-run global 0.6728에 \textbf{−0.038}, 같은-run XGB(full pool 12.07M) 0.7548에 −0.120, exp21 0.7411 비교선(방향만)에도 크게 미달. net correction \textbf{−1,926} (helpful 78,978 vs harmful 80,904, override precision 0.494 =\allowbreak{} 동전던지기), ρ=\allowbreak{}0.803(제안)/\allowbreak{}0.333(실행) — sparse가 아님. oracle gap recovery \textbf{−0.18}.

\textbf{그러나 Phase-1 질문(§26)의 답은 YES}: dense oracle macro \textbf{0.8833} (global +\allowbreak{}0.211) — brute 1.00 · inf 0.93 · dos 0.99 · ddos 1.00. 상보성은 실재하고 크다. 무너진 곳은 expert 구성도 oracle도 아니라 \textbf{advantage routing/\allowbreak{}verification 층}.

\textbf{메커니즘 분해} (1b/\allowbreak{}2a/\allowbreak{}3a/\allowbreak{}4b/\allowbreak{}5a):

\par\noindent \textbullet\enspace \textbf{잔차 블록의 benign 포화}: cluster 0이 1.84M(92\%, 전 클래스 혼합), cluster 1–3은 benign 슬리버  \ensuremath{\rightarrow}  expert 2–4의 block 186k가 \textbf{99.98\% benign} (2a). class-balanced weight는 클러스터링에는 들어가지만 TopB는 절대 점수 상위라서 \textbf{개수가 이긴다} — benign 1.59M의 중간 잔차가 rare의 고잔차를 수로 밀어냄. expert 1만 진짜 혼합 전문가(brute 57k·ddos 57k·bot 21k· dos 22k·inf 17k).


\par\noindent \textbullet\enspace \textbf{"balanced context vs natural context" 우위가 advantage를 오염}: anchor (클래스당 2k)가 모든 expert를 준-균형 context로 만들어, 1M 자연비율 global 대비 \textbf{공격 행 어디서나} balanced loss 우위(3a: 전 expert가 공격 클래스 frac\_\allowbreak{}positive \ensuremath{\approx} 0.9–1.0)  \ensuremath{\rightarrow}  positive-advantage 행 93\%(364,577/\allowbreak{}391,486), k* 타깃이 "regime 특화"가 아니라 "아무 expert나 호출"을 학습. λ\_\allowbreak{}cost=\allowbreak{}0이 이를 방치.


\par\noindent \textbullet\enspace \textbf{Δ 스케일 폭주로 verifier 붕괴}: w\_\allowbreak{}web=\allowbreak{}1133  \ensuremath{\rightarrow}  Δ 평균이 web \textasciitilde{}500–680 vs benign \textasciitilde{}0.01 (3a). pseudo-Huber(slope 1)가 사실상 L1이 되고 holdout pearson r이 expert 2·4에서 \textbf{음수}(−0.20/\allowbreak{}−0.48), MAE 176(4b). 수락된 1.34M override 중 절반이 유해 — brute .72 \ensuremath{\rightarrow} .40, inf .42 \ensuremath{\rightarrow} .14 붕괴(benign-포화 expert 2–4가 공격 행을 잠식). 반면 ddos .91 \ensuremath{\rightarrow} .97, dos .51 \ensuremath{\rightarrow} .74, web .16 \ensuremath{\rightarrow} .20, benign .98 \ensuremath{\rightarrow} .99는 개선 — 층이 제대로 서면 이득이 실재함을 방증.


\par\noindent \textbullet\enspace 비용: 총 49.2분(global fit 5.4m · eval global 12.3m · dense 10.7m) — dedup+\allowbreak{}cache로 exp21 저녁-단위 대비 대폭 절감. 평균 TFM 호출 1.80회.


\textbf{다음 원노브 후보} (각 1 run; 우선순위 제안):

\par\noindent \textbullet\enspace (a) \textbf{verifier를 pairwise ranking loss로} — PDF가 명시한 대안. Δ 절대 스케일(문제 3)을 우회하고 부호/\allowbreak{}순서만 학습.


\par\noindent \textbullet\enspace (b) \textbf{λ\_\allowbreak{}cost > 0} — 3a의 Δ 분포 기록이 근거: benign 평균 |Δ|\ensuremath{\approx}0\textasciitilde{}0.2 vs 공격 10–600  \ensuremath{\rightarrow}  λ·c\_\allowbreak{}k \ensuremath{\approx} 1이면 박빙-benign 제안 대량 제거, sparse 회복.


\par\noindent \textbullet\enspace (c) \textbf{S\_\allowbreak{}k의 per-class cap}(block benign 상한) — 문제 1 직접 타격, 단 TopB 정의 수정이라 스펙 변경 명시 필요.


\par\noindent \textbullet\enspace (d) γ<1 tempering (w\_\allowbreak{}web 1133 완화) — 문제 3 완화, Δ 정의 전체에 파급.


\subsubsection{§4 보강 — override 손익의 expert별 분해 (system\_\allowbreak{}dump 재분석, 사용자 검토 세션)}
\begingroup
\footnotesize
\begin{longtable}{>{\raggedright\arraybackslash}p{\dimexpr 0.24074\linewidth-2\tabcolsep\relax}>{\raggedleft\arraybackslash}p{\dimexpr 0.14815\linewidth-2\tabcolsep\relax}>{\raggedleft\arraybackslash}p{\dimexpr 0.14815\linewidth-2\tabcolsep\relax}>{\raggedleft\arraybackslash}p{\dimexpr 0.14815\linewidth-2\tabcolsep\relax}>{\raggedleft\arraybackslash}p{\dimexpr 0.16667\linewidth-2\tabcolsep\relax}>{\raggedright\arraybackslash}p{\dimexpr 0.14815\linewidth-2\tabcolsep\relax}}
\toprule
expert & 수락 & helpful & harmful & precision & net \\ \midrule
\endfirsthead
expert & 수락 & helpful & harmful & precision & net \\ \midrule
\endhead
{}e1 (혼합 진짜 expert) & 229,724 & 56,927 & 112 & \textbf{0.998} & \textbf{+\allowbreak{}56,815} \\
{}e2 (benign 클론) & 343,631 & 13,644 & 78,809 & 0.148 & \textbf{−65,165} \\
{}e3 (benign 클론) & 249,076 & 4,070 & 112 & 0.973 & +\allowbreak{}3,958 \\
{}e4 (benign 클론) & 516,946 & 4,337 & 1,871 & 0.699 & +\allowbreak{}2,466 \\
\bottomrule
\end{longtable}
\endgroup
\textbf{전체 피해 80,904 중 96\%가 단일 confusion 쌍}: e2가 global이 맞힌 brute\_\allowbreak{}force 77,344행을 infiltration으로 뒤집음(186k benign 블록 +\allowbreak{} anchor 2k/\allowbreak{} 클래스 context에서 brute가 inf로 오독). 피해 성립에 3중 실패가 모두 필요: router가 e2–4를 구분 못해 brute 행을 e2로 오배정(원인: 클론 context  \ensuremath{\rightarrow}  동일 d\_\allowbreak{}k) × e2의 퇴화 posterior × verifier가 수락(Δ 회귀 붕괴). 반면 \textbf{e1은 precision 0.998·net +\allowbreak{}57k} — expert가 제대로 구성되면 advantage-routing 구조 자체는 이미 작동함(helpful 상위: benign \ensuremath{\rightarrow} brute 28.9k, benign \ensuremath{\rightarrow} ddos 26.8k 구조). 결정 증거 추가: cluster 3은 weighted mass 0.15\%·mean\_\allowbreak{}residual 0.49 (global이 잘 맞히는 행들)인데 centroid 1개를 차지 — 질량이 아니라 \textbf{거리가 centroid를 배정}(z-공간 heavy-tail 기하, τ=\allowbreak{}7,776).  \ensuremath{\rightarrow}  knob 우선순위 조정: 상류 φ/\allowbreak{}클러스터 층(TabPFN embedding 또는 rank 변환, 질량-집중 클러스터링)이 1순위, verifier ranking loss는 안전판 복구로 병행 후보.

\subsection{§5. EXP22B 사전 등록 — 원노브: φ 교체 (raw RobustScaler+\allowbreak{}PCA  \ensuremath{\rightarrow}  global-TabPFN embedding+\allowbreak{}PCA) (결과 전 고정)}
사용자 지시: "1번 phi 교체 코드 수정하고 run까지 해줘." 스크립트: \texttt{tabpfn/\allowbreak{}\allowbreak{}n\allowbreak{}fv3\_\allowbreak{}\allowbreak{}v3\_\allowbreak{}\allowbreak{}e\allowbreak{}xp22b\_\allowbreak{}\allowbreak{}ra\allowbreak{}cepfn\_\allowbreak{}\allowbreak{}ph\allowbreak{}i.\allowbreak{}py} (exp22 사본; 기존 파일 무수정). \textbf{원노브}: φ =\allowbreak{} RobustScaler \ensuremath{\rightarrow} PCA16 (raw 46 feat)  \ensuremath{\rightarrow}  \textbf{global TabPFN get\_\allowbreak{}embeddings(192-d)  \ensuremath{\rightarrow}  PCA16} (\texttt{-\allowbreak{}-\allowbreak{}phi-\allowbreak{}mo\allowbreak{}de embed\allowbreak{}\_\allowbreak{}\allowbreak{}pca}; raw\_\allowbreak{}pca=\allowbreak{}exp22 재현, quantile\_\allowbreak{}pca=\allowbreak{}저비용 대안 옵션 동봉). PDF §6.2의 "이후 TabPFN embedding과 비교" 단계. 구조적 파급 명시: ① φ가 global fit \textbf{이후} 구축(embedding이 global 필요)  \ensuremath{\rightarrow}  medoid context-selection 미지원(exp22 기록도 random이었음), ② h\_\allowbreak{}pre의 z가 "cheap"이 아니게 됨(get\_\allowbreak{}embeddings 추가 forward 1회/\allowbreak{}샘플; timings에 기록, exp21 선례로 1M-context embedding chunk 100k 상한). 그 외 pool·weight·K·TopB· anchor·advantage·router·verifier·결정 규칙 전부 exp22와 동일. 6b(override expert별 분해) 아티팩트와 CSV+\allowbreak{}PNG 자동 렌더링 추가(기록 형식 확장, 수치 무관).

\textbf{판정 기준(고정)}: 같은-run 튜플 vs exp22 (system 0.6346 /\allowbreak{} global 0.6728 /\allowbreak{} XGB 0.7548 /\allowbreak{} oracle 0.8833). 관심: ① 클러스터가 공격 regime을 분리하는가 — "mass 0.15\% 클러스터가 centroid 차지" 유형의 거리-납치 소멸 여부·τ 스케일(1b), ② expert 2–4 benign 클론 소멸(2a), ③ e2형 brute \ensuremath{\rightarrow} inf 대량 유해 override 소멸(6b), ④ \textbf{예상 잔존 균열(격리 실험 의도)}: positive-advantage \textasciitilde{}93\%는 anchor 구조 탓이므로 잔존 예상, verifier Δ 회귀 취약(γ=\allowbreak{}1 스케일)도 잔존 예상 — 이 run은 φ 층의 효과만 격리. 실행: capped 스모크  \ensuremath{\rightarrow}  uncapped cic2018 기록 run (--dense-eval) 자동 체인.

\subsubsection{§5 보강 — exp22b diff-리뷰 결과 및 수정 (기록 run 전 적용)}
3-렌즈 리뷰(원노브 순수성/\allowbreak{}임베딩 런타임/\allowbreak{}누수·순서): \textbf{누수·순서 PASS}, major 1건 =\allowbreak{} \textbf{해석 교란(수치 오염 아님)} — embed 모드에서 phi-fit 행 \textasciitilde{}17\%와 aff0 참조행 100\%가 C0 내부 행이라 "암기된 행의 test 임베딩"으로 PCA 기반/\allowbreak{} affinity 기준이 형성됨(자기 라벨 사본에 attend 가능한 regime). headline 지표는 비오염(관련 라벨 전부 D\_\allowbreak{}context; router/\allowbreak{}verifier 특징은 fit/\allowbreak{}eval 동일 생성). 조치: docstring에 confound 선언 +\allowbreak{} C0 오버랩 카운트를 로그·timings.json에 기록 (--phi-fit-exclude-c0 분리는 후속 knob 후보). minor 수정: quantile n\_\allowbreak{}quantiles=\allowbreak{}min(1000,n)·임베딩 로그 stage 태그·embed\_\allowbreak{}chunk 가드·концat 메모리 셰이브·phi ready에 실현 z\_\allowbreak{}dim 출력. 스모크(\texttt{20260825\allowbreak{}\_\allowbreak{}\allowbreak{}161943}, capped): τ 38,430 \ensuremath{\rightarrow} \textbf{89.1}, 클러스터 4개가 실제 공격 regime으로 분리(brute+\allowbreak{}dos/\allowbreak{}혼합 공격 질량 0.61/\allowbreak{}inf 경계/\allowbreak{}bot·ddos — benign 슬리버 소멸), override precision 0.9814, macro 0.8683. 수정 후 마이크로 재검증  \ensuremath{\rightarrow}  uncapped 기록 run.

\subsection{§6. EXP22B 기록 run 결과 (run \texttt{20260825\allowbreak{}\_\allowbreak{}\allowbreak{}174233}, uncapped cic2018, test 4.02M) — \textbf{φ 교체로 첫 global 상회 (+\allowbreak{}0.038), 그러나 새 퇴화 양식}}
\textbf{판정}: racepfn\_\allowbreak{}system macro \textbf{0.7084} — 같은-run global 0.6700 \textbf{+\allowbreak{}0.0384 (최초 상회)}, net correction \textbf{+\allowbreak{}53,562} (helpful 57,177 vs harmful 3,615, precision \textbf{0.9405}), ρ=\allowbreak{}0.229·수락 0.034  \ensuremath{\rightarrow}  평균 TFM 호출 \textbf{1.23회 (sparse 회복)}. exp22(0.6346, net −1,926, ρ 0.80) 대비 전면 개선. 남은 거리: 같은-run XGB 0.7548 −0.046, exp21 0.7411(방향 비교) −0.033. oracle 0.8255(↓, 후술), gap recovery 0.247. 클래스: brute .7246 \ensuremath{\rightarrow} .8546 · dos .5312 \ensuremath{\rightarrow} .6999 · ddos .9407 \ensuremath{\rightarrow} .9603 · inf .3490 \ensuremath{\rightarrow} .3926 · benign .9812 \ensuremath{\rightarrow} .9898 · \textbf{web .1649 \ensuremath{\rightarrow} .0634 (악화, 후술)}. helpful 구조: benign \ensuremath{\rightarrow} brute 24.2k · benign \ensuremath{\rightarrow} ddos 12.1k · benign \ensuremath{\rightarrow} dos 10.7k · inf FP 정리 9.9k — 의도된 "global이 놓친 공격 회수" 그대로.

\textbf{φ 효과 확인 (1b)}: τ 7,776 \ensuremath{\rightarrow} 130.9. 거리-납치 소멸 — 질량 0.15\% 클러스터 없음(질량 13\textasciitilde{}48\%), 클러스터가 공격 구조를 가짐: c0=\allowbreak{}순수 공격(bot 20k·brute 40k·dos 20k, benign 18행!) · c1=\allowbreak{}benign/\allowbreak{}ddos 경계(ddos 132k 전량) · c2=\allowbreak{}benign+\allowbreak{}inf · c3=\allowbreak{}benign+\allowbreak{}inf/\allowbreak{}web.

\textbf{새 퇴화 — hard-attack 클론 (2a)}: 그런데 expert 4개의 block이 다시 거의 동일(brute 57,189 전량·dos \textasciitilde{}30k·inf \textasciitilde{}18k·web 252가 전원 포함; ddos만 684/\allowbreak{}1/\allowbreak{}0/\allowbreak{}0으로 차별화). 원인 산수: 클러스터 간 proximity 계수 exp(−d²/\allowbreak{}τ) =\allowbreak{} 0.02\textasciitilde{}0.41 (1.7자릿수)인데 \textbf{r은 benign 0.005  \ensuremath{\leftrightarrow}  web 1527로 6자릿수} — s\_\allowbreak{}ik =\allowbreak{} r·exp(−d²/\allowbreak{}τ)에서 r이 지배해 TopB가 클러스터 무관하게 전역 고-r 행으로 수렴. exp22의 benign-클론과 반대 방향의 같은 병: TopB의 r-대-근접성 스케일 불일치. ddos만 r이 작아(0.04) 근접성이 작동. 후속 원노브 후보: (c-1) S\_\allowbreak{}k를 hard assignment(assign=\allowbreak{}=\allowbreak{}k)로 제한 — PDF §7.2 "cluster k와 가깝고"의 의도에 더 충실, (c-2) --tau-c 축소(예: 중앙 클러스터간 d²의 1/\allowbreak{}10), (c-3) r을 클러스터 내 정규화.

\textbf{web 악화의 회계}: override가 web FP \textbf{7,640행} 생성 (true web 교체는 104행뿐). 대부분 global도 틀렸던 행(wrong \ensuremath{\rightarrow} wrong)이라 helpful/\allowbreak{}harmful 집계에 안 잡힘 — \textbf{net correction이 못 보는 클래스-정밀도 손상 경로}. w\_\allowbreak{}web=\allowbreak{}1133이 Δ에서 web 방향 교체를 과대평가하는 것과 일관(γ 완화/\allowbreak{}λ\_\allowbreak{}cost의 근거 추가).

\textbf{예상대로 잔존 (사전 등록 ④)}: verifier 여전히 붕괴(sign\_\allowbreak{}acc \textasciitilde{}0.25, pearson \textasciitilde{}0.1) — 수락률 14.7\%의 보수성 덕에 정밀도가 유지된 것이지 verifier가 좋아진 게 아님. e1은 84 수락 중 62 유해(음의 net) — 단 절대량이 작아 무해. router 타깃도 클론 탓에 e2로 쏠림(283k/\allowbreak{}300k).

\textbf{부수 관찰}: 같은 seed·같은 C0인데 global 지표가 run 간 이동(exp22 inf .4247·ddos .9073  \ensuremath{\leftrightarrow}  exp22b .3490·.9407) — CUDA 커널/\allowbreak{}할당 상태에 따른 비트-비결정성이 경계 행을 뒤집는 것으로 보임. 같은-run 비교 원칙(CLAUDE.md)이 이를 정확히 방어함; run 간 절대치 비교 금지 재확인. 비용: 총 71분 (eval\_\allowbreak{}router 754s =\allowbreak{} h\_\allowbreak{}pre 임베딩 — 선언된 "z 더는 cheap 아님" 실측).

다음: z-산점도(클래스/\allowbreak{}클러스터 색, 이 run의 μ·τ 재현) 생성 — \texttt{scripts/\allowbreak{}\allowbreak{}exp22b\_\allowbreak{}\allowbreak{}z\allowbreak{}\_\allowbreak{}\allowbreak{}scatter\allowbreak{}.\allowbreak{}py}, 산출물은 run 디렉토리에 7a로 저장.

\subsubsection{§6 보강 — z-공간 산점도 (\texttt{7a\_\allowbreak{}\allowbreak{}z\_\allowbreak{}\allowbreak{}sca\allowbreak{}tter.\allowbreak{}png\allowbreak{}}, scripts/\allowbreak{}exp22b\_\allowbreak{}z\_\allowbreak{}scatter.py)}
D\_\allowbreak{}mine 클래스-캡 표본 139,610행을 run과 동일 seed로 φ 재현 후 z[0:2] 투영 (배정은 run의 μ·τ 사용; 표본 median d² 158.9 \ensuremath{\approx} run τ 130.9 — 기하 재현 확인). 판독: ① 공간이 유계 연결 매니폴드(좌표 \textasciitilde{}±18 — raw φ의 5.5만 거리 소멸), ② ddos+\allowbreak{}bot이 고립 섬, dos·brute가 각각 밀집 영역, ③ \textbf{infiltration이 benign 구름 내부/\allowbreak{}인접에 상주} — inf가 어려운 이유와 c2/\allowbreak{}c3가 benign+\allowbreak{}inf 혼합인 이유가 시각적으로 확정, ④ 클러스터는 매니폴드를 연속 영역으로 분할(μ 4개 모두 매니폴드 내부 — 거리-납치 없음), ⑤ 클론 원인도 보임: 고-r 공격 영역들이 모든 μ에서 동일 자릿수 거리(d² 10²\textasciitilde{}10²·5)라 exp(−d²/\allowbreak{}τ)가 r을 못 이김.



\section*{출처와 실행 추적}

상세 수치의 원시 실행은 \path{inventory/runs.csv}, 실험별 연결은 \path{inventory/experiments.csv}에서 확인한다. 표·그림의 재현 입력은 \path{evidence/}와 \path{data/}에 둔다. 문서 이동·통합 내역은 \path{configs/experiments/migrations/20261020_documents/manifest.json}에 기록했다.

\end{document}
